\section{Discussion}
\label{sec:discussion}

\placeholder{Opening paragraph: restate the main result once B7 and B8 are in.}

\paragraph{What scales and what does not.} Exact inference across the boundary works,
and it is calibrated, but it stops at M0. The cost of an exact posterior grows with
the number of genes and the length of the record, and the full reduced model is
already above our budget. The approximate route removes the stochastic paths from
the sampler. Two quantities limit it. The first is the Jensen gap, which we declare
and correct (\cref{sec:jensen}). The second is $R$, the variance of the bulk
log-likelihood over the paths that the transcripts allow. $R$ is about
\pn{info.Rmzero} on M0 and \pn{info.Rcorea} on the full reduced model at 10\% assay
noise. So the metabolites constrain the paths together, and a single draw of paths
per parameter value is not enough. The Gaussian bulk factor handles this regime.
\placeholder{Confirm once B7 reports.}

\paragraph{The observation model sets the interface.} The three graphs of
\cref{sec:seam} give a general rule. Single-cell data inform the parameters of their
own module. Information between modules travels on bulk observables, and it follows
the gain of the coupling. A modeller can compute the gains before any inference, and
in this work the gains predicted the information in each direction before any
posterior existed.

\paragraph{Reductions as approximations to inference.} A dynamical reduction changes
the simulator and keeps full Bayesian inference. A statistical cut keeps the
simulator and changes the target. We prefer the first, because its cost can be
measured against a more detailed composition. The formalism swap is such a reduction.
We do not claim that the more detailed composition is the truth about the cell. We
claim that each choice of description is declared, and that its cost is measured.

\subsection{Limitations}

\begin{itemize}
  \item \emph{Synthetic data only.} The data come from the model itself. Real
  assays need observation models for unknown cell numbers in bulk metabolomics and
  for capture efficiency in single-cell RNA counts.
  \item \emph{One organism, one reduction.} We demonstrate the framework on one
  reduced model. \placeholder{O3: build the generic registry, or narrow the
  framework claim to ``demonstrated on one organism''.}
  \item \emph{Asserted priors.} Many priors have no width in their source, so we set
  the width. The uncertainty budget (\cref{sec:budget}) reports how much this choice
  matters. \placeholder{Once built.}
  \item \emph{Exact calibration at small scale only.} Coverage of the exact sampler
  at the full reduced model is out of reach at about \pn{cost.fullscale} CPU-hours per
  posterior.
  \item \emph{Weak reverse coupling.} The approximate route relies on the small
  elasticities of the rebuild (\pn{gain.txlo}--\pn{gain.tl}). A perturbation that
  drains the nucleotide pools toward the half-saturation constant would remove this
  basis.
\end{itemize}

\subsection{Future work}

\placeholder{One short paragraph, in this order: (1) the method's own limits:
condition the bulk factor on the observed transcripts; path sampling only for pools
near the one-particle floor; strong-feedback regimes. (2) The full well-stirred
\syn{} model through the syn3a-ensemble port, then real data. (3) Comparing model
structure and formalisms, which the swap makes possible, by predictive scores rather
than Bayes factors. (4) Goal-directed inference and experimental design: which parts
of a coupled simulator must be kept, and which latent variables inferred, to keep
the information a prediction needs?}
